Anh Diệu như tôi được biết - Nguyễn Ngọc Giao

Nguồn: \href{https://www.diendan.org/nhung-con-nguoi/anh-dieu-nhu-toi-duoc-biet}{Báo Diễn Đàn}

Có những người mà ta thường nghĩ ngay tới mỗi khi có một sự việc, sự kiện, tình huống mới hay phức tạp xảy ra, để tự hỏi : không biết anh/chị ấy nghĩ gì? Phan Đình Diệu, với tôi, là một người như vậy. Cũng lạ, vì anh sống ở Hà Nội, trong bầu không khí nhiều năm khá ngột ngạt, tôi ở Paris, chỉ có thể liên lạc với nhau bằng điện thoại hay qua mạng, đều là những kênh truyền thông mà người ta có thể nghe/đọc trộm, nên không bao giờ chúng tôi trao đổi trực tuyến những vấn đề "nhạy cảm" . Còn nhớ trong những năm 90, thời gian mà anh có những bản kiến nghị gửi giới lãnh đạo, hay những bài phát biểu trong những cuộc họp của Mặt trận Tổ quốc, mỗi khi nhận được văn bản (qua những con đường khác nhau), mệt nhất đối với ban biên tập Diễn Đàn là kiểm tra tính xác thực (có phải của anh Diệu không, có câu chữ nào bị sửa đổi, xuyên tạc không) trước khi công bố. Chúng tôi không thể hỏi thẳng tác giả, càng không thể "xin phép đăng", vì sẽ tạo cớ cho nhà cầm quyền quy kết anh "phát tán tài liệu nội bộ" (một vụ tương tự là lá thư của ông Võ Văn Kiệt gửi Bộ chính trị ĐCSVN, chúng tôi nhận được 2 văn bản chụp, một bản có chữ viết tay mà chúng tôi nhận ra nét chữ của tác giả, nhưng không biết người gửi là ai, sau khi Diễn Đàn công bố, lá thư được phổ biến khá rộng rãi ở nước ngoài, thì trong những tội trạng mà toà án "nhân dân" Hà Nội quy kết để giam tù các anh Nguyễn Kiến Giang và Lê Hồng Hà, có tội "gửi cho báo Diễn Đàn").

Thành ra những dịp duy nhất tôi có thể trao đổi thẳng thắn với anh là những lần được gặp anh. Không kể lần đầu gặp nhau, trong một tình huống đặc biệt : tháng 9 năm 1970, tôi về Hà Nội (lần đầu sau năm 1954, di cư), trong phái đoàn Việt kiều tại Pháp được mời về dự Quốc khánh 2-9. Lần ấy, tôi dở dang luận án về giải tích phức, định xoay ngang làm toán ngôn ngữ (với tham vọng là dùng lí thuyết phân tích chuỗi của Harris để làm mô hình cấu trúc câu tiếng Việt). Các anh Uỷ ban Khoa học Kỹ thuật xếp cho tôi một giờ gặp anh Diệu để anh hướng dẫn. Gần nửa thế kỉ đã trôi qua, tôi không nhớ anh đã giảng cho tôi những gì về ngữ pháp tạo sinh của Chomsky và tình hình "dịch thuật tự động" hồi đó. Tham vọng của tôi cũng sẽ dở dang vì vấp phải trở ngại mà cho đến nay ngôn ngữ học tiếng Việt vẫn chưa vượt qua : thế nào là một từ (mot/word). Điều chắc chắn là năm ấy, chúng tôi chưa hề trao đổi với nhau về thực trạng đất nước.

Phải đợi sau năm 1975, anh có thời gian sang Pháp ở mấy tháng, chúng tôi mới có dịp gặp nhau và nói chuyện thoải mái, dù là gặp riêng, hay gặp nhiều người. Từ đó cho đến cách đây hơn mười năm, thỉnh thoảng anh có dịp qua Pháp, tôi đều được gặp anh. Nhất là trong nửa đầu thập niên 2000, vợ chồng Hà Dương (con gái anh, nối nghiệp cha, cũng như em trai là Dương Hiệu) ở cũng xóm Les Juilliottes (Maisons Alfort) với tôi. Thành ra chúng tôi chủ yếu gặp gỡ chuyện trò ở Paris. Vài lần về nước (tôi bị cấm cửa khá lâu, từ 1982 đến 2001, 2003-2005, và từ 2008 đến nay), tôi đều lại phố Đội Cấn thăm anh Diệu chị Hương, nhưng không thuận tiện để trao đổi sâu. Lần duy nhất chúng tôi gặp nhau ngoài Paris và Việt Nam là cuộc hội thảo "Vietnam Update" (tháng 12-1994) của trường Đại học Quốc gia Úc (Canberra). Lần đó, ban tổ chức có mời cả anh Lữ Phương, nhưng anh Phương không được xuất ngoại. Tôi thì đi nửa vòng trái đất bằng một cái giấy "thông hành" (laissez-passer) xếp gọn như một cái thẻ cư trú, nhưng dàn ra thì dài loằng ngoằng như một cái đàn accordéon : giờ chót sở công an Pháp cấp vì tôi, tuy không có quy chế Hansen (vô tổ quốc), nhưng không còn hộ chiếu Việt Nam. Anh Diệu, sau mấy năm bị "cấm vận" vì những bài phát biểu, dường như chính quyền đã thấy hố khi báo chí quốc tế (đặc biệt là tuần báo Far Eastern \& Economic Review) điểm mặt chỉ tên, nên đã quyết định "bỏ cấm vận". Nếu tôi nhớ không lầm, trong khoảng thời gian 1991-1992, tổng bí thư Đỗ Mười đã gặp anh và lắng nghe đề nghị của anh là tổ chức một hội nghị trí thức Việt Nam (trong nước và ngoài nước) để trao đổi về tình hình đất nước. Anh Diệu có nhắn tôi viết thư cho ông Đỗ Mười để hoan nghênh dự án này. Sau đó, không nhớ ở Canberra năm 1994 hay sau đó ở Paris, anh kể ông Đỗ Mười hứa sau khi đi thăm Trung Quốc về, sẽ chính thức trả lời. Nhưng rồi, câu trả lời chưa hề tới tai anh. Còn lá thư của tôi, không được trả lời, nhưng có nhận được thư "biên nhận" của Văn phòng Tổng bí thư. Nói tới việc này, không phải để khoe (có gì đáng khoe), mà để ghi nhận một cung cách làm việc mà rất tiếc chính quyền Việt Nam không có (và đến nay vẫn chưa có) thói quen.

Tôi là một nhà báo nghiệp dư (tuy cũng đã 60 năm viết lách) hạng tồi : tôi không có thói quen ghi chép trong hay sau các cuộc gặp. Một phần là tôi không thích phỏng vấn chính thức, mà chỉ ưa trao đổi "off", nên không ghi chép cũng để người đối thoại nói năng thoải mái hơn. Phần nữa, tôi cậy có một trí nhớ khá tốt, đến khi tuổi cao, bắt đầu quên, mới học đến chữ ngờ.

Nay anh Diệu đã ra đi, nếu muốn ghi lại trung thực một số ý kiến của anh, hay một số tình tiết liên quan tới anh, sẽ phải bỏ nhiều thời gian đọc lại các bài viết của anh và những bài cũ của tôi liên quan tới anh. Thời gian không cho phép, may sao chiều nay được đọc bài của Huy Đức trên trang Facebook (Diễn Đàn vừa đăng lại ở đây). Tôi có thể khẳng định là những sự việc mà anh Huy Đức kể lại về Phan Đình Diệu đều trung thực, phù hợp với những điều mà chính anh Diệu và hai ba người gần gũi anh đã kể lại cho tôi. Phải cảm ơn Huy Đức đã ghi chép cẩn thận, đúng tác phong của một nhà báo chuyên nghiệp cự phách. Nếu phải nói thêm, thì chỉ nhắc đến vai trò của Phan Đình Diệu trong việc thuyết phục chính quyền Việt Nam nối mạng internet năm 1997 ; sau này, các nhà nghiên cứu sẽ phải tìm hiểu sự đóng góp quyết định của ba nhân vật là Nguyễn Đình Ngọc, Chu Hảo và Phan Đình Diệu trong cuộc "cách mạng internet" ở Việt Nam. Và nếu phải nói tình tiết, tôi chỉ xin bổ sung trang FB của một tác giả khác mà tôi vừa đọc chiều nay nhưng quên không ghi. Tác giả có nói anh Diệu đã từng tuyên bố về "một lãnh đạo cấp cao nhất", đại ý : đồng chí ấy "rất vĩ đại", nhưng bây giờ kháng chiến đã thành công, đồng chí ấy sẽ "vĩ đại" hơn nữa nếu từ chức. Tôi có dịp hỏi và anh Diệu cười xác nhận : "lãnh đạo cao cấp nhất" mà hôm nay tác giả không tiện kể tên, anh Diệu nói rõ : đồng chí Lê Duẩn. Và câu ấy, anh phát biểu trong một buổi họp tổ ở Quốc hội. Đó là lúc anh làm đại biểu quốc hội, khoá thứ nhì – cũng là khoá cuối cùng. Còn chuyện quan hệ của Phan Đình Diệu với Đảng CSVN. Có một thời, không ít trí thức Việt Nam – tôi muốn nói những người trí thức chân chính – mong muốn đứng vào hàng ngũ của đảng để đóng góp vào việc nước. Đó cũng là thời người ta soi mói lí lịch đến ba bốn đời để không kết nạp, hay khi kết nạp nhỏ giọt thì để trưng bày tủ kính và vô hiệu hoá mọi ý kiến được cho là không "phải đạo". Khi tình hình xấu đi nhiều, thì "đảng ta" lại muốn kết nạp một số trí thức tiêu biểu, để chứng tỏ đã "đổi mới tư duy". Trong tình huống ấy, Phan Đình Diệu đã được đề nghị kết nạp. Và, theo anh kể, anh đã từ chối bằng câu trả lời ý nhị này : "tôi sẽ gia nhập đảng cộng sản khi nào có hai đảng cộng sản, và tôi sẽ chọn một trong hai".

Để kết thúc đôi dòng tưởng niệm hết sức riêng tư, cá nhân này, tôi chỉ xin nói gọn một ý về người anh vừa ra đi. Đối với tôi, Phan Đình Diệu là một người trí thức Việt Nam hiện đại theo nghĩa hoàn chỉnh nhất. Anh là một nhà khoa học say mê trọn đời với chuyên môn của mình, với môi trường cần thiết cho bộ môn tin học là công nghệ tin học mà Việt Nam phải xây dựng ; đồng thời, anh hết sức quan tâm tới những ngành khoa học xã hội và nhân văn, qua đó anh tìm hiểu, phân tích tình trạng xã hội, tính chất của chế độ kinh tế - chính trị Việt Nam, và từ đó, phát biểu ý kiến của mình với tư cách một công dân, một trí thức độc lập. Tư duy khoa học và độc lập, và phát biểu ý kiến một cách thẳng thắn và an nhiên nhất – cách thực hành hiệu quả nhất để giành lấy và bảo vệ quyền tự do ngôn luận – đó là tính cách của một "ông đồ Nghệ" với tất cả khí phách và thông tuệ cần thiết, nhưng là một "ông đồ Nghệ" hiện đại, đồng điệu với giới trí thức đúng nghĩa của mọi quốc gia hiện đại.

Nhiều người đã nhận xét đúng rằng xã hội Việt Nam mới có những phần tử trí thức, chưa có giới trí thức (intelligentsia). Bi kịch của Phan Đình Diệu, cũng là một bi kịch của nước ta. Là một trí thức chân chính, anh đã góp phần cao nhất để thành hình giới trí thức Việt Nam hiện đại, để xây dựng xã hội công dân Việt Nam. 

Nguyễn Ngọc Giao
